\section{Conclusions and immediate decision}
\label{sec:conclusions}

The apparent broad usefulness of DFPT should not be described as a mysterious
global cancellation.  DFPT computes a full, state-dependent static Kohn--Sham
response.  The missing many-body object is the phonon derivative of the
beyond-Kohn--Sham self-energy.  That derivative can be small for one perturbation
and large for another even when the equilibrium self-energy and mass enhancement
are large.

The clearest physical anchor is a uniform common potential coupled to conserved
total charge: it only changes the energy zero, so interactions cannot arbitrarily
renormalize it.  The Ward identity states this constraint precisely, but its
dynamic and static limits must not be confused.  The UEG extends the picture with
strong numerical evidence that a static density kernel remains accurate over
finite momenta up to $2k_F$ in the tested regime.  Because the UEG contains only
the density channel, this result suggests a compression principle rather than a
universal proof.

Real crystals add internal orbital, bond, hopping, hybridization, spin, and
interaction-modulation operators.  Existing DFT+DMFT, DFPT+$U$, and GWPT results
show that their corrections can be strongly mode-selective.  This motivates the
proposal's central hypothesis: the useful small parameter is not ``correlation
strength'' by itself, but the weight of a phonon in each operator channel times
the many-body amplification and dynamical variation of that channel.

The proposed next-generation calculation therefore retains DFPT, decomposes its
perturbation, and applies a channel kernel only where diagnostics demand it.  This
can be simpler and cheaper than computing a complete many-body vertex at every
mode and momentum while being more physically transparent than a black-box
correction to every matrix element.  It is not faster than the single DFPT
calculation used as its baseline.  Its current computational claim is narrower:
sparse higher-level anchors plus cheap inference can cost less than evaluating a
higher-level correction at every point.  A future speed claim against DFPT would
require a separately implemented surrogate for the DFPT coefficients and a
matched comparison with modern DFPT interpolation.

The accompanying software now fits a small channel-response matrix, predicts
held-out coefficient vectors, reports error, and compares DFPT-only, dense
higher-level, and sparsely corrected cost baselines.  It explicitly reports that
the present corrected path is not faster than DFPT.  Its bundled held-out vectors
are synthetic and its costs are normalized assumptions.  These results establish
reproducible software behavior, not material accuracy or a measured runtime
speedup.

\begin{keypoint}[First decision after submission]
The first research milestone should test one sharp statement: in SrVO$_3$, does a
basis-robust measure of traceless orbital-channel weight separate the strongly
enhanced Jahn--Teller mode from the weakly corrected breathing mode, while the UEG
density relation supplies the charge-channel control?  A positive result justifies
building the static prototype; a negative result forces the hypothesis to be
revised before software expansion.
\end{keypoint}
